% Section: conclusion (\input from main.tex)
\section*{Conclusion}
\addcontentsline{toc}{section}{Conclusion}

We implemented sequential A* on a 203\,396-node road graph of urban Ho Chi
Minh City built from OpenStreetMap, extended it to uncertain and dynamic
heuristics, to the lecturer's time-dependent speed model and to a
vector-valued $f$ (\NAMOA), implemented the four sorted-list concurrent sets
with three locks each, and combined both into parallel A*, alongside a locked
heap and \HDA. We also checked the lecture's pseudocode and reference notebook,
found ten problems in the notebook (the most serious makes every road segment
cost 1\,m) and provide a corrected, tested version. Every exact algorithm was
checked against an independent oracle for exact cost equality on thousands of
random queries, and the concurrent code also under ThreadSanitizer; all tests
pass on all four platforms. The study supports five conclusions:

\begin{enumerate}[nosep]
  \item \textbf{A better heuristic beats more threads.} ALT landmarks make a
        query $12\times$ smaller than Dijkstra; no parallel version of a
        single query came close to that gain.
  \item \textbf{Inexact heuristics cost effort, not much quality.} Uncertain
        and dynamic heuristics lose far less than their bound $1+\varepsilon$,
        but the re-openings they cause make the search slower. The lecturer's
        $h = d/\mathit{speed}(t)$ is not admissible and returns a slower route
        in most runs (87--91\%); dividing by the maximum speed fixes this.
  \item \textbf{The lazy list is the concurrent set that scales.} Its wait-free
        \texttt{contains} reaches $5\times$ on 10 cores. Fine-grained locking
        is the slowest design, and spinning without backoff collapses
        whenever threads outnumber cores.
  \item \textbf{A sorted list is a poor open list.} \texttt{removeMin} always
        contends at the head and insertion is linear. A list-based parallel A*
        is correct after four adaptations, but only the lazy list with an
        expensive heuristic beat sequential A* ($1.11\times$ on 4 performance cores on A; $1.55\times$ on 10 uniform Xeon cores on D).
  \item \textbf{Parallel A* needs coarse work per node.} On a road graph a node
        expansion takes less than a microsecond, and with a cheap heuristic
        every parallel version loses to sequential A*. With an expensive
        heuristic, \HDA{} and the shared heap reach $1.35$--$1.83\times$ ($1.91\times$ with a persistent pool and core pinning on D). For
        many short queries, running different queries in parallel is the
        better use of the cores.
\end{enumerate}
